\documentclass[11pt]{article}
\usepackage[utf8]{inputenc}
\usepackage[margin=1in]{geometry}
\usepackage{amsmath, amssymb, amsthm}
\usepackage{booktabs}
\usepackage{graphicx}
\usepackage{hyperref}
\usepackage{enumitem}
\usepackage{caption}

% Updated path to look 'up' one level to find the plots folder
\graphicspath{ {../plots/} }

\title{ML in Finance: Principal Component Analysis of Global Markets}
\author{Ed Hayes}
\date{March 2026}

\begin{document}

\maketitle

\section{Introduction}
Principal Component Analysis (PCA) is a foundational dimensionality reduction technique used in quantitative finance to identify latent drivers of asset returns[cite: 5]. Whether analyzing the term structure of interest rates or the sector rotation of the S\&P 500, PCA allows us to distill hundreds of correlated variables into a few interpretable factors[cite: 6].

\section{Mathematical Foundation}

\subsection{Data Representation and Stationarity}
Let $Y_t \in \mathbb{R}^n$ represent the vector of $n$ assets observed at time $t$[cite: 10]. To ensure the data is stationary, we transform the raw prices or yields[cite: 11]:
\begin{itemize}
    \item \textbf{Interest Rates:} We utilize \textit{first differences} ($\Delta y_t = y_t - y_{t-1}$) because yields are already expressed in percentages[cite: 12].
    \item \textbf{Equities:} We utilize \textit{percentage changes} or returns ($r_t = \frac{P_t - P_{t-1}}{P_{t-1}}$) to normalize for absolute price levels[cite: 13].
\end{itemize}

\subsection{The Covariance Matrix}
PCA seeks to find the directions of maximum variance[cite: 15]. Let $X$ be the $T \times n$ matrix of centered changes[cite: 15]. The sample covariance matrix $\Sigma$ is:
\[ \Sigma = \frac{1}{T-1} X^T X \] [cite: 17]
In financial applications, we typically perform PCA on the \textbf{Covariance Matrix} (unscaled) to preserve information about the relative volatility of different segments of the market[cite: 18].

\subsection{Eigen-Decomposition}
We diagonalize $\Sigma$ to find the principal components[cite: 20]:
\[ \Sigma = Q \Lambda Q^T \] [cite: 21]
where $\Lambda$ contains the eigenvalues (variance explained) and $Q = [v_1, v_2, v_3]$ contains the eigenvectors (factor loadings)[cite: 22].

\section{Cross-Asset Interpretation}
The first three principal components often explain over 95\% of market movement[cite: 25].

\begin{table}[h]
\centering
\begin{tabular}{@{}lll@{}}
\toprule
\textbf{Component} & \textbf{Fixed Income Interpretation} & \textbf{Equity Interpretation} \\ \midrule
PC1 & Level (Parallel Shift) & Market Factor (Beta / Systemic Risk) \\
PC2 & Slope (Twist) & Sector/Style Rotation (Growth vs. Value) \\
PC3 & Curvature (Butterfly) & Specific/Thematic Clusters \\ \bottomrule
\end{tabular}
\caption{Comparison of latent factors across Interest Rates and Equities[cite: 26, 27].}
\end{table}

\section{Historical Case Studies: 20-Year Regime Shifts}
By utilizing 5-year rolling windows, we observe how market regimes evolve[cite: 29]:

\subsection{The 2008 Global Financial Crisis (2005--2010)}
During this period, the Financials (XLF) sector showed a disproportionately high loading in PC1 and PC2[cite: 31]. In Fixed Income, the ``Slope'' factor became highly volatile as the Fed aggressively cut rates[cite: 32].

\subsection{The Pandemic and Inflation Era (2020--2026)}
In the current regime, Technology (XLK) has dominated the Market Factor[cite: 34]. In the Treasury market, the aggressive hiking cycle of 2022--2023 saw the ``Level'' factor regain nearly 90\% of explained variance[cite: 35].

\newpage
\appendix

\section{Interest Rate: Scree Plots (2000--2025)}
% Each image now on its own page
\begin{figure}[h!] \centering \includegraphics[width=0.95\textwidth]{Rates_2000_2005_Scree.png} \end{figure} \newpage
\begin{figure}[h!] \centering \includegraphics[width=0.95\textwidth]{Rates_2005_2010_Scree.png} \end{figure} \newpage
\begin{figure}[h!] \centering \includegraphics[width=0.95\textwidth]{Rates_2010_2015_Scree.png} \end{figure} \newpage
\begin{figure}[h!] \centering \includegraphics[width=0.95\textwidth]{Rates_2015_2020_Scree.png} \end{figure} \newpage
\begin{figure}[h!] \centering \includegraphics[width=0.95\textwidth]{Rates_2020_2025_Scree.png} \end{figure} \newpage

\section{Interest Rate: Factor Loadings (2000--2025)}
\begin{figure}[h!] \centering \includegraphics[width=0.95\textwidth]{Rates_2000_2005_Loadings.png} \end{figure} \newpage
\begin{figure}[h!] \centering \includegraphics[width=0.95\textwidth]{Rates_2005_2010_Loadings.png} \end{figure} \newpage
\begin{figure}[h!] \centering \includegraphics[width=0.95\textwidth]{Rates_2010_2015_Loadings.png} \end{figure} \newpage
\begin{figure}[h!] \centering \includegraphics[width=0.95\textwidth]{Rates_2015_2020_Loadings.png} \end{figure} \newpage
\begin{figure}[h!] \centering \includegraphics[width=0.95\textwidth]{Rates_2020_2025_Loadings.png} \end{figure} \newpage

\section{Equity Sector: Scree Plots (2005--2025)}
\begin{figure}[h!] \centering \includegraphics[width=0.95\textwidth]{Equity_2005_2010_Scree.png} \end{figure} \newpage
\begin{figure}[h!] \centering \includegraphics[width=0.95\textwidth]{Equity_2010_2015_Scree.png} \end{figure} \newpage
\begin{figure}[h!] \centering \includegraphics[width=0.95\textwidth]{Equity_2015_2020_Scree.png} \end{figure} \newpage
\begin{figure}[h!] \centering \includegraphics[width=0.95\textwidth]{Equity_2020_2025_Scree.png} \end{figure} \newpage

\section{Equity Sector: Factor Loadings (2005--2025)}
\begin{figure}[h!] \centering \includegraphics[width=0.95\textwidth]{Equity_2005_2010_Loadings.png} \end{figure} \newpage
\begin{figure}[h!] \centering \includegraphics[width=0.95\textwidth]{Equity_2010_2015_Loadings.png} \end{figure} \newpage
\begin{figure}[h!] \centering \includegraphics[width=0.95\textwidth]{Equity_2015_2020_Loadings.png} \end{figure} \newpage
\begin{figure}[h!] \centering \includegraphics[width=0.95\textwidth]{Equity_2020_2025_Loadings.png} \end{figure}

\end{document}